
\documentclass[10pt]{article}
\usepackage[T1]{fontenc}
\usepackage[utf8]{inputenc}
\usepackage{lmodern}
\usepackage{amsmath,amssymb,amsthm,mathtools}
\usepackage{tikz-cd}
\usepackage{geometry}
\usepackage{enumitem}
\geometry{margin=0.84in}
\setlength{\parindent}{1.25em}
\setlength{\parskip}{0.16em}
\emergencystretch=3em
\setlist{leftmargin=2.0em,itemsep=0.08em,topsep=0.18em}
\newcommand{\Acal}{\mathcal A}
\newcommand{\Bcal}{\mathcal B}
\newcommand{\Fcal}{\mathcal F}
\newcommand{\Gcal}{\mathcal G}
\newcommand{\Z}{\mathbb Z}
\newcommand{\Q}{\mathbb Q}
\newcommand{\muell}{\mu_{\ell^n}}
\newcommand{\Db}{D^b}
\newcommand{\Hom}{\operatorname{Hom}}
\newcommand{\Ext}{\operatorname{Ext}}
\newcommand{\Id}{\operatorname{Id}}
\newcommand{\Ker}{\operatorname{Ker}}
\newcommand{\pr}{\operatorname{pr}}
\newcommand{\R}{\mathrm R}
\newcommand{\colim}{\varinjlim}
\newcommand{\iso}{\xrightarrow{\sim}}
\newcommand{\og}{\guillemotleft}
\newcommand{\fg}{\guillemotright}
\newcommand{\toiso}{\xrightarrow{\sim}}

\providecommand{\C}{\mathbb C}
\providecommand{\Ql}{\mathbb Q_\ell}
\providecommand{\Zl}{\mathbb Z_\ell}
\providecommand{\N}{\mathbb N}
\providecommand{\Alb}{\operatorname{Alb}}
\providecommand{\Pic}{\operatorname{Pic}}
\providecommand{\Spec}{\operatorname{Spec}}
\providecommand{\Tr}{\operatorname{Tr}}


\providecommand{\og}{\guillemotleft}
\providecommand{\fg}{\guillemotright}


\providecommand{\D}{\mathrm D}

\providecommand{\lgop}{\operatorname{lg}}
\providecommand{\Ass}{\operatorname{Ass}}
\providecommand{\Ob}{\operatorname{Ob}}
\providecommand{\coker}{\operatorname{coker}}
\providecommand{\RHom}{\operatorname{RHom}}
\providecommand{\DR}{\operatorname{DR}}
\providecommand{\an}{\mathrm{an}}
\providecommand{\red}{\mathrm{red}}
\providecommand{\Gr}{\operatorname{Gr}}
\providecommand{\im}{\operatorname{im}}
\providecommand{\codim}{\operatorname{codim}}
\providecommand{\HH}{\mathbb H}
\providecommand{\Tot}{\operatorname{Tot}}
\providecommand{\End}{\operatorname{End}}
\providecommand{\rk}{\operatorname{rk}}
\providecommand{\Ocal}{\mathcal O}
\providecommand{\Ecal}{\mathcal E}
\providecommand{\Mcal}{\mathcal M}


\begin{document}

\begin{center}
{\Large\bfseries THÉORÈME DE LEFSCHETZ\\
ET CRITÈRES DE DÉGÉNÉRESCENCE\\
DE SUITES SPECTRALES}\\[0.6em]
par P. DELIGNE\footnote{Aspirant au F.N.R.S.}
\end{center}

\section*{1. Théorèmes généraux}

Soient $\Acal$ une catégorie abélienne et $\Db(\Acal)$ la sous-catégorie pleine de sa catégorie dérivée formée des complexes bornés. On rappelle que si $T$ est un foncteur cohomologique (voir [8] ou [9], p. 10) de $\Db(\Acal)$ dans une catégorie abélienne $\Bcal$, Verdier [8] a défini une suite spectrale, fonctorielle en $X\in\Db(\Acal)$:
\[
\tag{1.1}
E_2^{pq}=T(H^q(X)[p])\Longrightarrow T(X[p+q]).
\]

\paragraph{Proposition (1.2).} Soit $X\in\Db(\Acal)$; les conditions suivantes sont équivalentes:
\begin{enumerate}[label=(\roman*)]
\item quel que soit $T$ comme plus haut, la suite spectrale (1.1) dégénère;
\item l'objet $X$ est isomorphe, dans $\Db(\Acal)$, à $\sum_i H^i(X)[-i]$.
\end{enumerate}

\emph{(i) $\Rightarrow$ (ii).} Désignons par $T_i$ le foncteur cohomologique en $K$ défini par
\[
T_i(K)=\Hom_{\Db(\Acal)}(H^i(X),K).
\]
Pour chacun des $T_i$ ($i\in\Z$), la suite spectrale (1.1) s'écrit
\[
\tag{1.3}
E_2^{pq}=\Ext^p(H^i(X),H^q(K))\Longrightarrow
\Hom_{\Db(\Acal)}(H^i(X),K[p+q]),
\]
et l'homomorphisme de $\Hom(H^i(X),K[i])$ dans
\[
E_\infty^{0i}=\Hom(H^i(X),H^i(K))
\]
qui s'en déduit n'est autre que celui provenant de la fonctorialité de $H^0$.

Si les suites spectrales (1.3) dégénèrent, les flèches
\[
\Hom(H^i(X)[-i],K)\longrightarrow \Hom(H^i(X),H^i(K))
\]
sont surjectives. Par hypothèse, tel est le cas si $K=X$, de sorte qu'il existe des morphismes $a_i$ de $H^i(X)[-i]$ dans $X$, induisant l'identité sur le $i$-ième groupe de cohomologie. De plus, puisque $H^i(X)=0$ sauf pour un nombre fini de valeurs de $i$, les $a_i$ sont nuls sauf un nombre fini, ce qui permet de définir
\[
a=\sum_i a_i:\sum_i H^i(X)[-i]\longrightarrow X.
\]
Par construction, $a$ induit l'identité sur la cohomologie, donc est un isomorphisme.

\emph{(ii) $\Rightarrow$ (i).} La suite spectrale (1.1) est un foncteur additif en $X$, et dégénère lorsque $X$ n'a qu'un objet de cohomologie non nul, donc aussi si $X$ est somme de complexes n'ayant qu'un objet de cohomologie non nul.

\paragraph{Remarque (1.4).} Il résulte de la démonstration précédente qu'il suffit de vérifier la condition (i) pour les foncteurs cohomologiques du type $\Hom_{\Db(\Acal)}(A,*)$ où $A\in\operatorname{Ob}(\Acal)$. Par dualité, il suffirait aussi de considérer les seuls foncteurs $\Hom_{\Db(\Acal)}(*,A)$, la condition (ii) étant autoduale.

\paragraph{Théorème (1.5).} Soient $X\in\Db(\Acal)$, $n$ un entier et $u$ un endomorphisme de degré $2$ de $X$,
\[
u:X\longrightarrow X[2],
\]
dont les itérés induisent des isomorphismes
\[
u^i:H^{n-i}(X)\iso H^{n+i}(X)\qquad(i\ge0).
\]
Sous ces conditions, $X$ est isomorphe à $\sum_i H^i(X)[-i]$, dans la catégorie $\Db(\Acal)$.

Vérifions la condition (i) de (1.2). Soit donc $T$ un foncteur cohomologique de $\Db(\Acal)$ dans $\Bcal$, et posons $T^p(K)=T(K[p])$.

On appelle partie primitive de $H^{n-i}(X)$, et on désignera par ${}_0H^{n-i}(X)$, le noyau de l'homomorphisme
\[
u^{i+1}:H^{n-i}(X)\longrightarrow H^{n+i+2}(X).
\]
On vérifie que, pour $i\ge0$, les applications
\[
\tag{1.6}
\begin{gathered}
\bigoplus_{k\ge0} u^k:\bigoplus_{k\ge0}{}_0H^{n-i-2k}(X)\iso H^{n-i}(X),\\
\bigoplus_{k\ge0}u^{k+i}:\bigoplus_{k\ge0}{}_0H^{n-i-2k}(X)\iso H^{n+i}(X)
\end{gathered}
\]
sont des isomorphismes.

La suite spectrale (1.1) étant fonctorielle en $X$, et « compatible » aux translations sur les degrés, l'endomorphisme $u$ de $X$ définit un endomorphisme de degré $2$ de la suite spectrale
\[
u:E_r^{pq}\longrightarrow E_r^{p,q+2},
\]
commutant aux différentielles $d_r$ ($r\ge2$).

Lorsque $r=2$, l'endomorphisme $u$ de la suite spectrale
\[
u:T^p(H^q(X))\longrightarrow T^p(H^{q+2}(X))
\]
se déduit du morphisme $u:H^q(X)\to H^{q+2}(X)$ par fonctorialité de $T^p$.

Soit $i\ge0$. Désignons par ${}_0E_r^{p,n-i}$ le noyau de l'homomorphisme
\[
u^{i+1}:E_r^{p,n-i}\longrightarrow E_r^{p,n+i+2}
\]
(« partie primitive » de la suite spectrale). Pour $r=2$, les décompositions (1.6) induisent par fonctorialité de chaque $T^p$ des décompositions
\[
\tag{1.7}
\begin{gathered}
\bigoplus_{k\ge0}u^k:\bigoplus_{k\ge0}{}_0E_2^{p,n-i-2k}\iso E_2^{p,n-i},\\
\bigoplus_{k\ge0}u^{k+i}:\bigoplus_{k\ge0}{}_0E_2^{p,n-i-2k}\iso E_2^{p,n+i}.
\end{gathered}
\]

Prouvons par récurrence sur $r$ que les différentielles $d_r$ ($r\ge2$) sont nulles. L'hypothèse de récurrence implique que $E_2=E_r$, de sorte que la décomposition (1.7) s'applique au terme $E_r$ de la suite spectrale. Puisque $u$ commute aux $d_r$, il suffira de prouver que $d_r$ s'annule sur la partie primitive de $E_r$.

Le diagramme suivant est commutatif:
\[
\begin{tikzcd}[column sep=3.2em,row sep=2.6em]
{}_0E_r^{p,n-i} \arrow[r,"d_r"] \arrow[d,"u^{i+1}"'] & E_r^{p+r,n-i-r+1} \arrow[d,"u^{i+1}"]\\
E_r^{p,n+i+2} \arrow[r,"d_r"'] & E_r^{p+r,n+i-r+3}.
\end{tikzcd}
\]
La flèche verticale gauche est nulle par définition de la primitivité. La flèche verticale droite est un monomorphisme, car la flèche
\[
u^{i+r-1}:E_r^{p+r,n-(i+r-1)}\longrightarrow E_r^{p+r,n+(i+r-1)}
\]
est un isomorphisme et $r-1\ge1$. Il en résulte que $d_r$ est nul.

\paragraph{Remarque (1.8).} On prendra garde qu'on ne peut pas en général choisir un isomorphisme entre $X$ et $\sum_iH^i(X)[-i]$ qui transforme $u$ en la somme des applications $H^i(X)\to H^{i+2}(X)$. On obtient aisément un contre-exemple en prenant $H^i(X)=0$ pour $i\ne n$.

On peut définir un isomorphisme « plus beau que les autres » entre $X$ et $\sum_iH^i(X)[-i]$, mais on n'aura pas à s'en servir ici.

\paragraph{Remarque (1.9).} Soient $X\in\Db(\Acal)$, $n$ et $s$ des entiers et $u$ un endomorphisme de degré $s$ (resp. $2s$) de $X$. On suppose que $H^i(X)=0$ lorsque $i$ n'est pas de la forme $ks$ ($0\le k\le n$) (resp. $0\le k\le2n$) et que $u^{n-2k}$ (resp. $u^i$) induit un isomorphisme entre $H^{ks}(X)$ et $H^{(n-k)s}(X)$ (resp. entre $H^{(n-i)s}(X)$ et $H^{(n+i)s}(X)$). Sous ces hypothèses, la démonstration de (1.5) montre encore que $X$ est isomorphe à $\sum_iH^i(X)[-i]$.

\paragraph{Remarque (1.10).} Soit $(X_i)_{i\in\Z}$ une famille d'objets de $\Db(\Acal)$ et $u$ une famille d'homomorphismes de degré $2$: $u_i:X_i\to X_{i+1}[2]$. Si, quels que soient $i$ et $j$, $u_i$ induit un isomorphisme entre $H^{n-j}(X_i)$ et $H^{n+j}(X_{i+j})$, il est encore vrai que chaque $X_i$ est isomorphe à la somme de ses objets de cohomologie; pour le voir, il suffit d'appliquer formellement les raisonnements qui précèdent à
\[
u:\sum_iX_i\longrightarrow\sum_iX_i,
\]
même si cette somme n'existe pas dans $\Db(\Acal)$.

La même remarque s'applique à la variante (1.9).

\paragraph{Théorème (1.11).} Soit $X\in\Db(\Acal)$ et supposons que $X$ admette des endomorphismes de degré $0$, $\pi_i:X\to X$, tels que $H^j(\pi_i)=\delta_{ij}$. Alors $X$ est isomorphe à $\sum_iH^i(X)[-i]$ dans la catégorie $\Db(\Acal)$.

Appliquons encore le critère (1.2) (i). Le diagramme
\[
\begin{tikzcd}[column sep=3.0em,row sep=2.2em]
E_r^{p,q} \arrow[r,"d_r"] \arrow[d,"\pi_q"'] & E_r^{p+r,q-r+1} \arrow[d,"\pi_q"]\\
E_r^{p,q} \arrow[r,"d_r"'] & E_r^{p+r,q-r+1}
\end{tikzcd}
\]
est commutatif, la première colonne est l'identité, et la seconde zéro, de sorte que $d_r=0$.

\paragraph{Corollaire (1.12).} Soit $X\in\Db(\Acal)$ la somme d'une famille finie d'objets $X_k\in\Db(\Acal)$. Si on a
\[
X\simeq\sum_iH^i(X)[-i],
\]
alors, pour tout $k$, on a
\[
X_k\simeq\sum_iH^i(X_k)[-i].
\]
Soient $j_k$ l'injection canonique de $X_k$ dans $X$, $\pr_k$ la projection canonique de $X$ sur $X_k$ et $\pi_i$ un endomorphisme de $X$ tel que $H^j(\pi_i)=\delta_{ij}$. Pour chaque $k$, les $\pr_k\pi_i j_k$ vérifient l'hypothèse de (1.11).

\paragraph{Remarque (1.13).} On vérifie que, pour que l'isomorphisme entre $X$ et $\sum_iH^i(X)[-i]$ puisse être choisi tel que les $\pi_i$ s'identifient aux $\pr_i$, il faut et il suffit que les $\pi_i$ soient des idempotents deux à deux orthogonaux. Il existe alors un seul isomorphisme ayant cette propriété et induisant l'identité sur la cohomologie.

\section*{2. Applications}

Soit $f:X\to Y$ une application continue entre espaces topologiques, voire un morphisme de topos, $A$ un faisceau d'anneaux sur $X$ et $\Fcal$ un faisceau de $A$-modules. Si $u\in H^2(X,A)$, $u$ définit un endomorphisme de degré $2$ dans $D^b(X)$, encore noté
\[
u:\Fcal\longrightarrow\Fcal[2].
\]
Par fonctorialité, $u$ définit un endomorphisme de degré $2$:
\[
Rf_*(u):Rf_*\Fcal\longrightarrow Rf_*\Fcal[2].
\]
On dira que $(X,f,u,\Fcal,n)$ vérifie la condition de Lefschetz, ou que $\Fcal$ vérifie la condition de Lefschetz relativement à $u$, si, pour chaque $i\ge0$, les flèches obtenues en itérant $i$ fois la flèche $Rf_*(u)$,
\[
Rf_*(u)^i:R^{n-i}f_*\Fcal\longrightarrow R^{n+i}f_*\Fcal\qquad(i\ge0),
\]
sont des isomorphismes. Le théorème (1.5) donne ici:

\paragraph{Proposition (2.1).} Si $\Fcal$ vérifie la condition de Lefschetz relativement à $u$, alors, dans $D^b(Y)$, on a
\[
Rf_*\Fcal\simeq\sum_iR^if_*\Fcal[-i]
\]
et en particulier la suite spectrale de Leray
\[
H^p(Y,R^qf_*\Fcal)\Longrightarrow H^{p+q}(X,\Fcal)
\]
dégénère.

Si on y avait tenu, on aurait pu supposer que $\Fcal$ soit un $(A,f^*B)$-bimodule, où $B$ est un faisceau d'anneaux sur $Y$, et énoncer (2.1) dans $D^b(Y,B)$.

\paragraph{(2.2)} Le cas des schémas, munis de la topologie étale, ne rentre pas dans le cadre (2.1). On supposera pour la suite que la théorie des $\Z_\ell$-faisceaux et $\Q_\ell$-faisceaux (non nécessairement constructibles) ait été faite, et que la catégorie dérivée de la catégorie de ces « faisceaux » soit raisonnable. En ce qui concerne les énoncés de dégénérescence de suites spectrales, le lecteur qui n'aurait pas confiance en le travail futur de Jouanolou pourra vérifier, en revenant à la démonstration de (1.5), qu'on peut s'en passer. Soit $\ell$ un nombre premier. On suppose pour toute la suite que $\ell$ est inversible sur tous les schémas considérés.

Rappelons que, sur tout schéma $X$, le $\Z_\ell$-faisceau $\Z_\ell(1)$ est défini par la formule
\[
\Z_\ell(1)=\varprojlim_n\mu_{\ell^n}.
\]
Ce faisceau est un $\Z_\ell$-module inversible, ce qui permet de définir, pour tout $\Z_\ell$-faisceau ou $\Q_\ell$-faisceau $\Fcal$, le « faisceau » $\Fcal(i)$ par la formule
\[
\Fcal(i)=\Fcal\otimes_{\Z_\ell}\Z_\ell(1)^{\otimes i}\qquad(i\in\Z).
\]
On aura, pour tout morphisme $f:X\to Y$ de schémas et tout couple de $\Z_\ell$- ou $\Q_\ell$-faisceaux $\Fcal$ et $\Gcal$ sur $X$ et $Y$ respectivement:
\[
\tag{2.3}
Rf_*(\Fcal)(i)=Rf_*(\Fcal(i))\qquad\text{et}\qquad f^*(\Gcal(i))=f^*(\Gcal)(i).
\]
Si $\mathcal L$ est un faisceau inversible sur $X$, sa première classe de Chern $\ell$-adique se trouve dans $H^2(X,\Z_\ell(1))$.

Si $f:X\to Y$ est un morphisme de schémas, si $\Fcal$ est un $\Z_\ell$-faisceau, respectivement $\Q_\ell$-faisceau, sur $X$ et si $u\in H^2(X,\Z_\ell(1))$, respectivement $u\in H^2(X,\Q_\ell(1))$, on dira que $\Fcal$ vérifie la condition de Lefschetz relativement à $u$ si, pour chaque $i\ge0$, les flèches obtenues en itérant $i$ fois la flèche $Rf_*(u)$,
\[
Rf_*(u)^i:R^{n-i}f_*\Fcal\longrightarrow R^{n+i}f_*\Fcal(i)\qquad(i\ge0),
\]
sont des isomorphismes. En vertu de (2.3), cela implique que toutes les flèches
\[
Rf_*(u)^i:R^{n-i}f_*\Fcal(k)\longrightarrow R^{n+i}f_*\Fcal(k+i)
\]
sont des isomorphismes. Appliquant (1.10), on trouve:

\paragraph{Proposition (2.4).} Avec les notations précédentes, si $\Fcal$ vérifie la condition de Lefschetz relativement à $u$, alors, dans $D^b(Y,\Z_\ell)$, respectivement $D^b(Y,\Q_\ell)$, on a
\[
\tag{2.5}
Rf_*\Fcal\simeq\sum_iR^if_*\Fcal[-i]
\]
et en particulier la suite spectrale de Leray
\[
H^p(Y,R^qf_*\Fcal)\Longrightarrow H^{p+q}(X,\Fcal)
\]
dégénère.



Si $g$ est un morphisme de schémas de $Y$ dans $S$, (2.4) implique encore la dégénérescence de la suite spectrale de Leray
\[
R^p g_* R^q f_*\mathcal F\Longrightarrow R^{p+q}(gf)_*\mathcal F.
\]

\paragraph{(2.6.1)} Reste à passer en revue quelques cas intéressants où la condition de Lefschetz est vérifiée; nous donnerons des détails justificatifs dans (2.7). Soient donnés $f:X\to Y$, $u$, $\mathcal F$ et $n$. Dans le cas topologique, on sait d'après Godement [2], II, (4.7.1), que si $f$ est propre et $Y$ localement paracompact, la formation des images directes supérieures commute au passage aux fibres, de sorte que la condition de Lefschetz se vérifie fibre par fibre. Dans le cadre des schémas, la même réduction s'applique, la référence à Godement étant remplacée par SGA 4, XII, (5.1). Dans le cadre des schémas encore, si $f$ est propre et lisse et si $\mathcal F$ est constant tordu, c'est-à-dire lisse dans une terminologie nouvelle, les $R^if_*\mathcal F$ sont encore constants tordus (cf. SGA 4, XVI, (2.2)); de sorte que, si $Y$ est connexe, il suffit même de vérifier la condition de Lefschetz sur une fibre géométrique.

Dans tous les exemples donnés, $f$ sera une application continue, respectivement un morphisme de schémas, propre et lisse de dimension relative topologique $2n$, respectivement algébrique $n$. Une application continue est dite lisse si, localement à la source, elle est isomorphe à la projection de $\R^n\times Y$ sur $Y$.

\paragraph{(2.6.2)} Si $f$ est un morphisme propre et lisse de variétés kählériennes, le faisceau constant $\C$ vérifie la condition de Lefschetz relativement à la $2$-classe fondamentale de $X$.

\paragraph{(2.6.3)} Si $X$ est un schéma relatif complexe (voir [3]), projectif et lisse sur l'espace topologique localement paracompact $Y$, et muni d'un faisceau inversible relativement ample $\mathcal O(1)$, le faisceau constant $\C$ vérifie la condition de Lefschetz relativement à la première classe de Chern de $\mathcal O(1)$.

\paragraph{(2.6.4)} Supposons que $f$ soit un morphisme de schémas projectif et lisse, avec $Y$ connexe, et que soit donné un faisceau inversible relativement ample $\mathcal O(1)$ sur $X$. Si l'on peut remonter en caractéristique $0$ une fibre géométrique de $f$, et sa polarisation, alors le faisceau $\Ql$-adique constant $\Ql$ vérifie la condition de Lefschetz relativement à la première classe de Chern $\ell$-adique de $\mathcal O(1)$.

On sait qu'on peut relever en caractéristique $0$, au sens précédent, dans les cas suivants:
\begin{itemize}
\item $X$ est une variété de Severi--Brauer sur $Y$;
\item $X$ est une intersection complète non singulière d'hypersurfaces dans une variété de Severi--Brauer sur $Y$, pour la polarisation induite par la polarisation canonique de la variété de Severi--Brauer;
\item $X$ est un schéma abélien polarisé sur $Y$ (Mumford, à paraître).
\end{itemize}
On conjecture que (2.6.4) reste vrai lorsqu'on ne sait pas se remonter en caractéristique $0$, mais ce n'est prouvé que pour les surfaces.

\paragraph{(2.6.5)} Si $f$ est un morphisme projectif et lisse de dimension relative $2$ et si $X$ est muni d'un faisceau inversible relativement ample $\mathcal O(1)$, alors le faisceau constant $\Ql$ vérifie la condition de Lefschetz relativement à la première classe de Chern de $\mathcal O(1)$.

\paragraph{(2.7)} Dans les cas (2.6.2) et (2.6.3), la $2$-classe de cohomologie donnée sur $X$ induit sur les fibres de $f$ la $2$-classe définie par la structure kählérienne des fibres. Ainsi, pour prouver (2.6.2) et (2.6.3), il suffit, d'après (2.6.1), de prouver (2.6.2) dans le cas particulier où $Y$ est réduit à un point. L'assertion résulte alors de la théorie de Hodge et de la décomposition de Hodge--Lepage (Weil [10], IV, n\textsuperscript{o} 6, corollaire au théorème 5). La démonstration donnée dans loc. cit. s'applique encore si $\mathcal F$ est un faisceau localement constant de $\C$-vectoriels de dimension finie, muni d'une structure hermitienne définie positive localement constante.

D'après (2.6.1), pour vérifier (2.6.4), il suffit de traiter le cas où $Y$ est le spectre d'un corps de caractéristique $0$, et le principe de Lefschetz permet de supposer $Y=\Spec(\C)$.

On sait alors (cf. SGA 4, XI, (4.4)) que, pour tout faisceau $\Ql$-adique constant tordu $\mathcal F$ sur $X$, si $\mathcal F^{an}$ désigne le faisceau localement constant de $\Ql$-vectoriels sur l'espace topologique $X^{an}$ déduit de $\mathcal F$, alors les $\Ql$-vectoriels de cohomologie, algébrique et transcendante, de $\mathcal F$ et $\mathcal F^{an}$, sont isomorphes:
\[
\tag{2.7.2}
H^i(X,\mathcal F)\iso H^i(X^{an},\mathcal F^{an}).
\]
Choisissant un plongement de $\Ql$ dans $\C$, on en déduit des isomorphismes
\[
\tag{2.7.3}
H^i(X,\mathcal F)\otimes_{\Ql}\C\iso
H^i(X^{an},\mathcal F^{an})\otimes_{\Ql}\C.
\]
Sur $\C$, l'application exponentielle définit un isomorphisme dit canonique entre $\Z/\ell^n$ et $\mu_{\ell^n}$, donc entre $\Zl$ et $\Zl(1)$. L'image, par l'isomorphisme composé
\[
H^2(X,\Zl(1))\longrightarrow H^2(X,\Zl)\longrightarrow H^2(X^{an},\Z)\otimes\Zl,
\]
de la première classe de Chern $\ell$-adique de $\mathcal O(1)$ coïncide avec la première classe de Chern transcendante de $\mathcal O(1)$. L'isomorphisme (2.7.3) ramène donc le cas auquel on s'était réduit de (2.6.4) à (2.6.3).

Pour la démonstration de (2.6.5), voir Kleiman [6], pp. 28 et suivantes. Le point essentiel est le théorème suivant (Weil [11], p. 127).

\paragraph{(2.8)} Soient $S$ une surface projective lisse sur un corps $k$ et $D$ une courbe lisse tracée sur $S$ telle que $\mathcal O(D)$ soit ample. Alors l'application composée
\[
\Pic^0(S)\longrightarrow\Pic^0(D)=\Alb(D)\longrightarrow\Alb(S)
\]
est une isogénie.

\paragraph{Remarque (2.9).} Le théorème de dégénérescence de suite spectrale (2.2) déduit de (2.6.4) a été conjecturé par Grothendieck, par des considérations de \og poids\fg, lorsque $Y$ est projectif et lisse sur un corps algébriquement clos.

\paragraph{Remarque (2.10).} Serre m'a fait remarquer que les théorèmes de dégénérescence de suites spectrales déduits de (2.6.2) et (2.6.3) peuvent aussi se démontrer par une extension facile de la méthode utilisée par Blanchard ([1], II, 1), tout au moins si l'on dispose de la dualité de Poincaré sur la base. Réciproquement, la méthode suivie ici permet de compléter les théorèmes II (1.1) et II (1.2) de loc. cit., la démonstration de l'implication (ii) $\Rightarrow$ (i) étant essentiellement celle de Blanchard, qui devait supposer les faisceaux images directes supérieures constants.

\paragraph{Théorème (2.11).} Soient $f:X\to Y$ une application continue propre et lisse (2.6.1) d'espaces topologiques de dimension relative $n$, avec $Y$ localement paracompact, $k$ un corps commutatif de caractéristique zéro et
$u\in H^0(Y,R^2f_*k)$. On suppose que:
\begin{enumerate}[label=\alph*)]
\item le faisceau $R^1f_*k$ est constant, c'est-à-dire simple;
\item les cup-produits itérés
\[
 u^i\wedge:R^{n-i}f_*k\longrightarrow R^{n+i}f_*k
\]
sont des isomorphismes.
\end{enumerate}
Alors les conditions suivantes sont équivalentes:
\begin{enumerate}[label=(\roman*)]
\item la classe $u$ est induite par un élément de $H^2(X,k)$;
\item la transgression
\[
d_2:H^0(Y,R^2f_*k)\longrightarrow H^2(Y,f_*k)
\]
est nulle sur $u$;
\item dans $\D^b(Y,k)$, on a
\[
Rf_*k\simeq\bigoplus_iR^if_*k[-i].
\]
\end{enumerate}
L'assertion (i) $\Rightarrow$ (iii) est contenue dans (2.1). L'assertion (iii) implique la dégénérescence de toute la suite spectrale de Leray
\[
\tag{2.12}
H^p(Y,R^qf_*k)\Longrightarrow H^{p+q}(X,k),
\]
et a fortiori (ii); il reste à prouver que (ii) $\Rightarrow$ (i).

Soit $x\in Y$ et $X_x=f^{-1}(x)$. La classe induite $u_x\in H^2(X_x,k)$ ne s'annule sur aucune composante connexe de la variété $X_x$, qui est donc orientable. Si l'on désigne par $\Tr_u$ l'application composée
\[
H^{2n}(X_x,k)\xrightarrow{(u_x^n)^{-1}}H^0(X_x,k)=H_0(X_x,k)\xrightarrow{\Tr}k,
\]
la dualité de Poincaré implique que la forme bilinéaire
\[
H^{n-i}(X_x,k)\times H^{n+i}(X_x,k)\longrightarrow k,
\qquad (x,y)\longmapsto \Tr_u(xy)
\]
est une dualité parfaite. Cette construction se globalise et fournit
\[
\Tr_u:R^{2n}f_*k\longrightarrow k.
\]

Prouvons tout d'abord que, dans (2.12), $d_2u\in H^2(Y,R^1f_*k)$ est nul. Si $x\in H^0(Y,R^1f_*k)$, pour une raison de degré, on a $xu^n=0$, donc
\[
d_2(xu^n)=d_2x\,u^n-nxu^{n-1}d_2u=0.
\]
D'après b), l'application
\[
 u^{n-1}\wedge:H^0(Y,R^1f_*k)\longrightarrow H^0(Y,R^{2n-1}f_*k)
\]
est un isomorphisme. Ainsi, si (ii) est vérifiée, quel que soit $y\in H^0(Y,R^{2n-1}f_*k)$, on a
\[
y\,d_2u=0,
\]
et en particulier
\[
\tag{2.13}
\Tr_u(y\,d_2u)=0.
\]
Soit $V$ un $k$-vectoriel tel que $R^1f_*k$ soit isomorphe au faisceau constant $V$. On a $d_2u\in H^2(Y,V)$ et, d'après (2.13), pour toute forme linéaire $w$ sur $V$, l'image de $d_2u$ par $w$ dans $H^2(Y,k)$ est nulle. On conclut que $d_2u=0$.

Si $d_2u=0$, la démonstration de (1.5) montre que toutes les différentielles $d_2$ de (2.12) sont nulles, de sorte que $E_2=E_3$. Pour une raison de degré, $u^{n+1}=0$, de sorte que
\[
d_3u^{n+1}=(n+1)u^n d_3u=0.
\]
Puisque $d_3u\in H^3(Y,R^0f_*k)$, on déduit de b) que $d_3u=0$. Pour une raison de degré, on a $d_ru=0$ pour $r>3$, et $u$ provient donc d'un élément de $H^2(X,k)$.

\paragraph{(2.14)} Dans le cadre des schémas, le résultat précédent reste valable, mais n'est utilisable que dans le cas \og géométrique\fg, plus précisément lorsque $\Ql$ est isomorphe à $\Ql(1)$ sur $Y$.

\paragraph{(2.15)} Le corollaire (1.12) permet d'étendre certains des énoncés qui précèdent à certains faisceaux non constants.

\paragraph{Proposition (2.16).} Soient $g$ et $f$ deux applications continues, respectivement deux morphismes de schémas, composables
\[
X\xrightarrow{g}Y\xrightarrow{f}Z,
\]
et $\mathcal F$ un faisceau, respectivement un $\Zl$-faisceau ou $\Ql$-faisceau, sur $X$. Supposons que
\begin{enumerate}[label=\alph*)]
\item $Rg_*\mathcal F\in\D^b(Y)$ et $Rg_*\mathcal F\simeq\bigoplus_iR^ig_*\mathcal F[-i]$;
\item $R(fg)_*\mathcal F\in\D^b(Z)$ et $R(fg)_*\mathcal F\simeq\bigoplus_iR^i(fg)_*\mathcal F[-i]$.
\end{enumerate}
Alors, pour tout $k$, on a
\[
Rf_*(R^kg_*\mathcal F)\simeq\bigoplus_i R^if_*(R^kg_*\mathcal F)[-i].
\]
D'après b), le complexe
\[
R(fg)_*\mathcal F=Rf_*Rg_*\mathcal F\simeq
\bigoplus_kRf_*(R^kg_*\mathcal F[-k])
\]
est somme de ses objets de cohomologie; d'après (1.12), les $Rf_*R^kg_*\mathcal F$, qui, à un décalage près, en sont facteurs directs grâce à a), jouissent de la même propriété.

\paragraph{Proposition (2.17).} Soient $\mathcal E$ un fibré vectoriel localement libre sur un schéma $S$, $\mathcal F$ un faisceau $\ell$-adique constant tordu sur $\mathbb P(\mathcal E)$, et $f$ la projection de $\mathbb P(\mathcal E)$ sur $S$. Dans la catégorie dérivée, on a
\[
Rf_*\mathcal F\simeq\bigoplus_iR^if_*\mathcal F[-i].
\]
L'espace projectif est simplement connexe, et sa cohomologie est sans torsion (SGA 1, XI, 1.1 et SGA 4, XVI, 2.2). On a donc
\[
\mathcal F=f^*f_*\mathcal F,
\qquad
Rf_*\mathcal F=Rf_*\Zl\otimes f_*\mathcal F,
\]
formule qui permet de se ramener au cas où $\mathcal F=\Zl$. Soit $u$ la classe de Chern du faisceau inversible $\mathcal O(1)$ sur $\mathbb P(\mathcal E)$. On sait alors que $\Zl$ vérifie la condition de Lefschetz relativement à $u$, et on conclut par (2.2).

\paragraph{(2.18)} De (2.17), on peut déduire, à l'aide de (2.16), le même résultat pour les fibrés en drapeaux de toutes espèces définis par $\mathcal E$. Dans le cadre topologique, ce qui précède s'applique tel quel aux fibrés vectoriels sur $\C$, et s'étend aux fibrés vectoriels sur $\mathbb H$ grâce à (1.9). Pour les fibrés réels, il faut se limiter au cas où $\mathcal F$ provient d'un faisceau de $2$-torsion sur la base.

\paragraph{Remarque (2.19).} Soit $X$ un schéma abélien de dimension relative $g$ sur une base $Y$. Pour tout entier $n$, la multiplication par $n$, soit $[n]$, définit un endomorphisme
\[
[n]^*:Rf_*\Ql\longrightarrow Rf_*\Ql.
\]
Prenant des combinaisons linéaires à coefficients rationnels, et à dénominateurs bornés en termes de $g$, de ces endomorphismes, on construit des endomorphismes
\[
\pi_i:Rf_*\Ql\longrightarrow Rf_*\Ql
\]
tels que $H^j(\pi_i)=\delta_{ij}$. Appliquant (1.11), on conclut, sans hypothèse projective, que
\[
Rf_*\Ql\simeq\bigoplus_iR^if_*\Ql[-i].
\]
L'idée exploitée ici est due à Lieberman.

\section*{3. La forme infinitésimale du théorème de semi-continuité}

Je rappelle dans ce paragraphe quelques résultats qu'on parvient à trouver dans EGA III, § 7. La formulation (3.4) m'a été signalée par L. Illusie.

\paragraph{(3.0)} Dans tout ce paragraphe, on désigne par $A$ un anneau local noethérien fixé une fois pour toutes, par $\mathfrak m$ son idéal maximal et par $k$ son corps résiduel. On désigne par $D_f^-(A)$ la sous-catégorie de la catégorie dérivée de la catégorie des $A$-modules, formée des complexes bornés supérieurement de $A$-modules de type fini.

\paragraph{Proposition (3.1).} Soit $M$ un module de type fini sur $A$.
\begin{enumerate}[label=(\roman*)]
\item Pour tout homomorphisme, pas nécessairement local, de $A$ dans un anneau artinien $B$, et tout $B$-module de type fini $N$, on a
\[
\tag{3.1.1}
\operatorname{lg}_B(M\otimes_A N)\ge
\operatorname{lg}_B(N)\cdot\dim_k(M\otimes_A k).
\]
\item Les conditions suivantes sur $M$ sont équivalentes:
\begin{enumerate}[label=\alph*)]
\item $M$ est un module libre;
\item quels que soient $B$ et $N$ comme en (i), on a
\[
\tag{3.1.2}
\operatorname{lg}_B(M\otimes_A N)=
\operatorname{lg}_B(N)\cdot\dim_k(M\otimes_A k);
\]
\item quel que soit $n\in\N$, la condition (3.1.2) est vérifiée pour $B=N=A/\mathfrak m^n$;
\item quel que soit $\mathfrak p\in\operatorname{Ass}(A)$, $M_\mathfrak p$ est plat sur $A_\mathfrak p$ et la condition (3.1.2) est vérifiée pour $B=N=A_\mathfrak p/\mathfrak pA_\mathfrak p$.
\end{enumerate}
\end{enumerate}



Si $A$ est sans composantes immergées, ces conditions équivalent encore à:
\begin{enumerate}[label=\alph*),start=5]
\item quel que soit le point maximal $\mathfrak p$ de $\Spec(A)$, on a
\[
\tag{3.1.2'}
\lgop_{A_\mathfrak p}(M_\mathfrak p)=\lgop(A_\mathfrak p)\cdot
\dim_k(M\otimes_A k).
\]
\end{enumerate}
La démonstration, une application du lemme de Nakayama, est laissée au lecteur.

\paragraph{(3.2)} Pour tout complexe $K$ et tout entier $n$, on désignera par $\tau_{\ge n}(K)$ le complexe défini par
\[
\tag{3.2.1}
\tau_{\ge n}(K)^i=
\begin{cases}
0,& i<n,\\
\coker(d^{n-1}),& i=n,\\
K^i,& i>n,
\end{cases}
\]
de sorte que
\[
\tag{3.2.2}
H^i(\tau_{\ge n}(K))=
\begin{cases}
0,& i<n,\\
H^i(K),& i\ge n.
\end{cases}
\]
Soit $K$ un objet de $D^-_{\mathrm{coh}}(A)$ (cf. (3.0)). Quitte à remplacer $K$ par un complexe isomorphe dans la catégorie dérivée, on peut supposer les composantes de $K$ libres de type fini.

Quels que soient $B$ et $N$ comme en (3.1) (i), on a, au niveau des complexes,
\[
\tau_{\ge n}(K\otimes_A^L N)\simeq \tau_{\ge n}(K)\otimes_A N.
\]
Prenant les caractéristiques d'Euler--Poincaré des deux membres, on trouve
\[
\sum_{i\ge0}(-1)^i\lgop_B\bigl(H^{n+i}(K\otimes_A^L N)\bigr)
=\lgop_B(\coker(d^{n-1})\otimes_A N)+
\sum_{i>0}(-1)^i\lgop_B(K^{n+i}\otimes_A N).
\]
Soustrayons de cette identité l'identité analogue pour $B=N=k$, multipliée par $\lgop_B(N)$. D'après (3.1) (ii), $a\Leftrightarrow b$, on a, les $K^i$ étant libres,
\[
\tag{3.2.3}
\begin{aligned}
&\sum_{i\ge0}(-1)^i\lgop_B\bigl(H^{n+i}(K\otimes_A^L N)\bigr)
-\lgop_B(N)\sum_{i\ge0}(-1)^i\dim_k\bigl(H^{n+i}(K\otimes_A^L k)\bigr)\\
&\qquad=\lgop_B(\coker(d^{n-1})\otimes_A N)
-\lgop_B(N)\dim_k(\coker(d^{n-1})\otimes_A k),
\end{aligned}
\]
ce qui permet d'appliquer (3.1) au premier membre de (3.2.3).

\paragraph{Théorème (3.3) (théorèmes d'échange et de semi-continuité).} Soient $K\in\Ob D^-_{\mathrm{coh}}(A)$ (cf. (3.0)) et $n\in\Z$.
\begin{enumerate}[label=(\roman*)]
\item Pour tout homomorphisme de $A$ dans un anneau artinien $B$ et tout $B$-module de type fini $N$, on a
\[
\tag{3.3.1}
\sum_{i\ge0}(-1)^i\lgop_B\bigl(H^{n+i}(K\otimes_A^L N)\bigr)
\le
\lgop_B(N)\sum_{i\ge0}(-1)^i\dim_k\bigl(H^{n+i}(K\otimes_A^L k)\bigr).
\]
\item Les conditions suivantes sur $K$ et $n$ sont équivalentes:
\begin{enumerate}[label=\alph*)]
\item Le foncteur cohomologique en le $A$-module $N$, $H^*(K\otimes_A^L N)$, vérifie la propriété d'échange en $*=n$, c'est-à-dire que les conditions équivalentes suivantes sur $K$ et $n$ sont vérifiées:
\begin{enumerate}[label=a.\arabic*)]
\item le foncteur $H^n(K\otimes_A^L N)$ est exact à gauche;
\item il existe un module $Q$ tel que le foncteur précédent soit isomorphe au foncteur $\Hom_A(Q,N)$;
\item le foncteur $H^{n-1}(K\otimes_A^L N)$ est exact à droite;
\item il existe un module $P$ tel que le foncteur précédent soit isomorphe au foncteur $P\otimes_A N$;
\item la flèche $H^{n-1}(K)\to H^{n-1}(K\otimes_A^L k)$ est surjective;
\item il existe un complexe borné supérieurement $K'$, quasi-isomorphe à $K$ et à composantes plates, respectivement libres de type fini, tel que $d'^{\,n-1}=0$.
\end{enumerate}
\item Quels que soient $B$ et $N$ comme en (i), on a
\[
\tag{3.3.2}
\sum_{i\ge0}(-1)^i\lgop_B\bigl(H^{n+i}(K\otimes_A^L N)\bigr)
=
\lgop_B(N)\sum_{i\ge0}(-1)^i\dim_k\bigl(H^{n+i}(K\otimes_A^L k)\bigr).
\]
\end{enumerate}
Si $A$ est sans composantes immergées, ces conditions équivalent encore à:
\begin{enumerate}[label=\alph*),start=3]
\item quel que soit le point maximal $\mathfrak p$ de $\Spec(A)$, on a
\[
\tag{3.3.2'}
\sum_{i\ge0}(-1)^i\lgop_{A_\mathfrak p} H^{n+i}(K_\mathfrak p)=
\lgop(A_\mathfrak p)\sum_{i\ge0}(-1)^i\dim_k\bigl(H^{n+i}(K\otimes_A^L k)\bigr).
\]
\end{enumerate}
\end{enumerate}

\paragraph{Remarque (3.3.3).} On eût pu allonger la liste par les conditions analogues à (3.1) (ii) c) et d).

L'assertion (i) résulte de (3.2.3) et de (3.1) (i). D'après EGA III, (7.4.2), et sa démonstration, les conditions a.1), a.3) et a.6) sont équivalentes et, si $K$ est représenté par un complexe borné supérieurement de modules libres de type fini, elles signifient encore que $\coker(d^{n-1})$ est libre; compte tenu de (3.1) (ii) $a\Leftrightarrow b\Leftrightarrow e$ et de (3.2.3), ceci prouve l'équivalence de a.1), a.3), a.6), b) et c). On a, trivialement, a.6) $\Rightarrow$ a.2) $\Rightarrow$ a.1) et a.6) $\Rightarrow$ a.4) $\Rightarrow$ a.3). L'équivalence de la condition a.5) avec les autres est contenue dans EGA III, (7.5.2).

\paragraph{Corollaire (3.4).} Soient $A$, $K$ et $n$ comme en (3.3), et $k\in\N$.
\begin{enumerate}[label=(\roman*)]
\item Quels que soient $B$ et $N$ comme en (3.3) (i), on a
\[
\tag{3.4.1}
\sum_{0\le i\le 2k}(-1)^i\lgop_B\bigl(H^{n+i}(K\otimes_A^L N)\bigr)
\le
\lgop_B(N)\sum_{0\le i\le 2k}(-1)^i\dim_k\bigl(H^{n+i}(K\otimes_A^L k)\bigr).
\]
\item Les conditions suivantes sur $K$, $n$ et $k$ sont équivalentes:
\begin{enumerate}[label=\alph*)]
\item Le foncteur cohomologique en $N$, $H^*(K\otimes_A^L N)$, vérifie la propriété d'échange en $*=n$ et en $*=n+2k+1$.
\item[\textup{a'}] Le complexe $K$ est quasi-isomorphe à un complexe borné supérieurement $K'$, à composantes plates, respectivement libres de type fini, tel que $d'^{\,n-1}=d'^{\,n+2k}=0$.
\item Quels que soient $B$ et $N$ comme en (i), on a égalité dans (3.4.1).
\end{enumerate}
Si $A$ est sans composantes immergées, ces conditions équivalent encore à:
\begin{enumerate}[label=\alph*),start=3]
\item quel que soit le point maximal $\mathfrak p$ de $\Spec(A)$, on a égalité dans (3.4.1) pour $B=N=A_\mathfrak p$.
\end{enumerate}
\end{enumerate}

\paragraph{(3.4.2)} Pour $k=0$, il s'agit là de propriétés du seul hypertor $H^n(K\otimes_A^L N)$, considéré comme foncteur en le $A$-module $N$. Elles s'expriment encore en disant que $H^n(K)$ est libre et que, pour tout $N$, la flèche canonique de $H^n(K)\otimes_A N$ dans $H^n(K\otimes_A^L N)$ est un isomorphisme.

\paragraph{(3.4.3)} D'autre part, si $K$ est un complexe parfait, on déduit de (3.4), pour $n\ll0$, un résultat analogue à (3.3), la sommation sur $i\ge0$ dans (3.3.1, 2, 2') étant remplacée par une sommation sur $i\le0$.

L'assertion (i) et l'équivalence de a), b) et c) résultent de (3.3) et de l'identité
\[
\sum_{0\le i\le 2k}(-1)^i x_i=
\sum_{0\le i}(-1)^i x_i+
\sum_{0\le i}(-1)^i x_{i+2k+1}.
\]
Pour prouver a'), on applique successivement (3.3) (ii) $a\Leftrightarrow a.6)$ à $(K,n)$ et à $(\tau_{\ge n}(K),n+2k+1)$.

\paragraph{(3.5)} Lorsque $A$ est artinien, (3.4) donne, pour $k=0$, l'inégalité
\[
\tag{3.5.1}
\lgop_A H^n(K)\le \lgop(A)\dim_k\bigl(H^n(K\otimes_A^L k)\bigr),
\]
et, si on a égalité dans (3.5.1), alors la condition d'échange est vérifiée en $n$ et $n+1$, et $H^n(K)$ est un $A$-module libre.

\section*{4. Morphisme de Gysin en cohomologie de Hodge}

Le paragraphe est rédigé d'après des notes de A. Grothendieck.

\paragraph{(4.0)} Soit $f:X\to Y$ un morphisme propre entre schémas $X$ et $Y$ lisses purement de dimension relative $n$ et $m$, sur une base $S$. On suppose que $S$ est noethérien et admet un complexe dualisant pour pouvoir référer à [4] pour la définition de $\R f^!$; cette restriction est sans doute inutile. Enfin, posons $d=n-m$.

\paragraph{(4.1)} Toute forme différentielle relative sur $Y$ définit par image réciproque une forme différentielle relative sur $X$, d'où un morphisme
\[
\tag{4.1.1}
f^*\Omega_{Y/S}^q=Lf^*\Omega_{Y/S}^q\longrightarrow\Omega_{X/S}^q.
\]
D'après la théorie de la dualité pour les morphismes lisses et la transitivité de $\R f^!$, on a canoniquement
\[
\tag{4.1.2}
\R f^!\Omega_{Y/S}^m=\Omega_{X/S}^n[d].
\]
Appliquons à $\Omega_{Y/S}^q$ et $\Omega_{Y/S}^m$ la formule d'induction
\[
\R f_*\RHom(K,L)=\RHom(Lf^*K,\R f^!L),
\]
qui s'obtient aisément à partir de la définition de $\R f^!$ par bidualité (Hartshorne [4], chap. VII, § 3). Puisque
\[
\RHom(\Omega_{Y/S}^q,\Omega_{Y/S}^m)\simeq\Omega_{Y/S}^{m-q},
\]
on trouve
\[
\R f^!\Omega_{Y/S}^{m-q}=\RHom(Lf^*\Omega_{Y/S}^q,\Omega_{X/S}^n)[d].
\]
La flèche (4.1.1) définit donc par transposition une flèche
\[
\tag{4.1.3}
\Omega_{X/S}^{n-q}=\RHom(\Omega_{X/S}^q,\Omega_{X/S}^n)
\longrightarrow \R f^!\Omega_{Y/S}^{m-q}[d].
\]
D'après la théorie de la dualité, se donner une flèche (4.1.3) revient à se donner une flèche
\[
\tag{4.1.4}
\R f_*\Omega_{X/S}^{n-q}\longrightarrow \Omega_{Y/S}^{m-q}[-d].
\]

\paragraph{Définition (4.2).} Soient $X$, $Y$, $f$, $S$, $n$, $m$ et $d$ comme en (4.0). On appelle morphisme de Gysin la flèche (4.1.4)
\[
\Tr_f:\R f_*\Omega_{X/S}^q\longrightarrow\Omega_{Y/S}^{q-d}[-d]
\]
et les flèches telles que
\[
\Tr_f:H^p(X,\Omega_{X/S}^q)\longrightarrow H^{p-d}(Y,\Omega_{Y/S}^{q-d})
\]
qui s'en déduisent.

\paragraph{Proposition (4.3).} Soit $f:X\to Y$ un morphisme propre et birationnel de schémas lisses sur un corps $k$. Les flèches
\[
f^*:H^p(Y,\Omega_Y^q)\longrightarrow H^p(X,\Omega_X^q)
\]
sont injectives.

On se ramène à supposer $X$ et $Y$ purement de dimension $n$. Le morphisme composé
\[
\Tr_f\circ f^*:\Omega_Y^q\longrightarrow\R f_*f^*\Omega_Y^q\longrightarrow\R f_*\Omega_X^q\longrightarrow\Omega_Y^q
\]
est l'identité, car il coïncide avec l'identité sur un ouvert dense. Appliquant le foncteur $H^p(Y,\ )$, on trouve que le morphisme composé
\[
\Tr_f\circ f^*:H^p(Y,\Omega_Y^q)\longrightarrow H^p(X,\Omega_X^q)\longrightarrow H^p(Y,\Omega_Y^q)
\]
est l'identité, ce qui prouve (4.3).

\section*{5. Cohomologie de Hodge et de De Rham}

\paragraph{(5.1)} Soit $f:X\to S$ un morphisme lisse de schémas. Par définition, les faisceaux de cohomologie de De Rham relative de $X$ sur $S$ sont donnés par
\[
H_{\DR}^*(X/S)=\R^*f_*(\Omega_{X/S}^\bullet),
\]
où $\Omega_{X/S}^\bullet$, le complexe de De Rham relatif, est un complexe différentiel $S$-linéaire. On dispose d'une suite spectrale
\[
\tag{5.1.1}
E_1^{pq}=\R^qf_*\Omega_{X/S}^p\Longrightarrow H_{\DR}^{p+q}(X/S).
\]

\paragraph{(5.2)} Lorsque $S=\Spec(\C)$ et que $X$ est propre et lisse, on sait par GAGA ([7]) que la cohomologie de chacun des faisceaux cohérents $\Omega_X^p$, ainsi donc que l'hypercohomologie de $\Omega_X^\bullet$, sont les mêmes au sens algébrique, topologie de Zariski, ou transcendant, topologie usuelle de $X^{\an}$; en particulier, $\Omega_{X^{\an}}^\bullet$ étant une résolution du faisceau constant $\C$, on a
\[
H_{\DR}^n(X)=H^n(X^{\an},\C)
\]
et la conjugaison complexe sur le faisceau constant $\C$ induit une conjugaison complexe, de nature transcendante, sur $H_{\DR}^n(X)$. La suite spectrale (5.1.1) s'écrit ici
\[
\tag{5.2.1}
E_1^{pq}=H^q(X,\Omega_X^p)\Longrightarrow H_{\DR}^{p+q}(X);
\]
on désignera par $F^p(H^n(X))$ la filtration décroissante qu'elle définit sur son aboutissement. D'après Weil [10], chap. IV, n\textsuperscript{o} 4, th. 2, la suite spectrale (5.2.1) peut se calculer comme suite spectrale du complexe double $H^0(X^{\an},\Omega^{p,q})$ filtré par le premier degré, $\Omega^{p,q}$ désignant le faisceau des formes différentielles $C^\infty$ bihomogènes de bidegré $(p,q)$ (Weil [10], chap. II, n\textsuperscript{o} 1).

Supposons maintenant que $X$ soit projective, donc $X^{\an}$ une variété kählérienne. Les formes harmoniques forment alors un sous-complexe double (Weil [10], chap. II, n\textsuperscript{o} 6, cor. 1 au th. 2) du complexe double précédent, sur lequel $d'$ et $d''$ s'annulent. L'inclusion de ce sous-complexe induit un isomorphisme sur les termes $E_1$ des suites spectrales définies par ces doubles complexes, filtrés par le premier degré, comme on le voit en considérant l'opérateur $H$ (composante harmonique), rétraction du complexe double $H^0(X^{\an},\Omega^{p,q})$ sur le sous-complexe des formes harmoniques et vérifiant
\[
1-H=d''(2\delta'G)+(2\delta'G)d''
\]
(Weil [10], chap. IV, n\textsuperscript{o} 1: (I) et lemme 3; n\textsuperscript{o} 3: cor. 2 et chap. II, n\textsuperscript{o} 5, th. 2 (IX)).

Si on désigne par $H^{p,q}(X)$ l'image dans $H^{p+q}(X^{\an},\C)$ de l'espace des formes harmoniques de type $(p,q)$, on a donc
\[
\tag{5.2.2}
H^n(X,\C)=\bigoplus_{p+q=n}H^{p,q}(X),
\]
\[
\tag{5.2.3}
F^p(H^n(X))=\bigoplus_{i\ge p}H^{i,n-i}(X),
\]
\[
\tag{5.2.4}
\overline{H^{p,q}(X)}=H^{q,p}(X)
\quad\text{(conjugaison complexe)},
\]
et la suite spectrale (5.2.1) dégénère.

Revenant au cas général où $X$ est seulement supposé propre et lisse sur $\C$, on posera
\[
\tag{5.2.5}
H^{p,q}(X)=F^p(H^{p+q}(X))\cap\overline{F^q(H^{p+q}(X))},
\]
de sorte que
\[
\tag{5.2.6}
\overline{H^{p,q}(X)}=H^{q,p}(X).
\]
D'après (5.2.3), cette notation est compatible avec celle utilisée lorsque $X$ est projectif. On posera encore
\[
\tag{5.2.7}
h^{p,q}=\dim H^q(X,\Omega_X^p).
\]

\paragraph{Proposition (5.3).} Soit $X$ un schéma propre et lisse sur $\C$:
\begin{enumerate}[label=(\roman*)]
\item la suite spectrale (5.2.1) dégénère;
\item on a
\[
\tag{5.3.1}
F^p(H^n(X))=\bigoplus_{i\ge p}H^{i,n-i}(X),
\]
et en particulier
\[
\tag{5.3.2}
H^n(X)=\bigoplus_{p+q=n}H^{p,q}(X)
\quad\text{(somme directe)},
\]
\[
\tag{5.3.3}
h^{p,q}=h^{q,p}=\dim H^{p,q}(X).
\]
\end{enumerate}

On se ramène à supposer $X$ connexe. Soit $N$ sa dimension. D'après le lemme de Chow et la résolution des singularités (Hironaka [5]), il existe un morphisme projectif et birationnel $g:X'\to X$, avec $X'$ projectif et lisse sur $\C$. La suite spectrale
\[
\tag{5.3.4}
E_1^{pq}=H^q(X,\Omega_X^p)\Longrightarrow H_{\DR}^{p+q}(X)
\]
s'envoie, par image réciproque, dans la suite spectrale analogue
\[
\tag{5.3.5}
E_1'{}^{pq}=H^q(X',\Omega_{X'}^p)\Longrightarrow H_{\DR}^{p+q}(X'),
\]
qui, d'après (5.2), dégénère. D'après (4.3), le morphisme $g^*:E_1\to E_1'$ est injectif. On en conclut, par récurrence sur $r\ge1$, que $d_r=0$, que $E_r=E_{r+1}$ et que $g^*$ injecte $E_{r+1}$ dans $E'_{r+1}$:
\[
\begin{tikzcd}[column sep=4.0em,row sep=2.1em]
E_r^{pq} \arrow[r,"d_r"] \arrow[d,"g^*"'] & E_r^{p+r,q-r+1} \arrow[d,"g^*"]\\
E_r'{}^{pq} \arrow[r,"0"'] & E_r'{}^{p+r,q-r+1} .
\end{tikzcd}
\]
Puisque les suites spectrales (5.3.4) et (5.3.5) dégénèrent, et que $g^*$ est injectif sur leur terme initial, il l'est encore sur leur aboutissement. Des formules (5.2.3), (5.2.4) appliquées à $X'$ résulte que
\[
F^p(H^n(X'))\cap\overline{F^{n-p+1}(H^n(X'))}=\{0\};
\]
on en déduit la formule analogue
\[
\tag{5.3.6}
F^p(H^n(X))\cap\overline{F^{n-p+1}(H^n(X))}=\{0\},
\]
car $g^*$ envoie $F^p(H^n(X))$ dans $F^p(H^n(X'))$ et commute à la conjugaison complexe. En particulier, on a
\[
\sum_{i\ge p}h^{i,n-i}+\sum_{i\ge n-p+1}h^{i,n-i}\le \dim H^n(X)\le \sum_i h^{i,n-i},
\]
ce qu'on peut écrire
\[
\tag{5.3.7}
\sum_{i\ge p}h^{i,n-i}\le\sum_{i\le n-p}h^{i,n-i}.
\]
Par dualité de Serre, $h^{i,j}=h^{N-i,N-j}$, l'inégalité (5.3.7), pour $N-n$, implique l'inégalité opposée, de sorte que
\[
\dim F^p(H^n(X))+\dim \overline{F^{n-p+1}(H^n(X))}=\dim H^n(X).
\]
D'après (5.3.6), on a donc
\[
\tag{5.3.8}
H^n(X)=F^p(H^n(X))\oplus\overline{F^{n-p+1}(H^n(X))}.
\]
La décomposition (5.3.8) induit sur $F^{p-1}(H^n(X))$, qui contient l'un des facteurs, une décomposition
\[
F^{p-1}(H^n(X))=F^p(H^n(X))\oplus H^{p-1,n-p+1}(X),
\]
et la formule (5.3.1) s'en déduit par récurrence décroissante sur $p$; (5.3.2) et (5.3.3) résultent aussitôt de (5.3.1) et (5.2.6), ce qui achève la démonstration de (5.3).

\paragraph{Corollaire (5.4).} Sous les hypothèses de (5.3), on a:
\begin{enumerate}[label=(\roman*)]
\item Une classe de cohomologie complexe $a$ sur $X$ est dans $H^{p,q}(X)$ si et seulement si elle peut se représenter par une forme fermée $\alpha$ de type $(p,q)$.
\item Toute classe de cohomologie peut se représenter par une forme $\alpha$ vérifiant $d'\alpha=d''\alpha=0$.
\item Si la forme $\alpha$ vérifie $d'\alpha=d''\alpha=0$, les conditions suivantes sont équivalentes:
\begin{enumerate}[label=\alph*)]
\item $(\exists\beta)\ \alpha=d\beta$;
\item $(\exists\beta)\ \alpha=d'\beta$;
\item $(\exists\beta)\ \alpha=d''\beta$.
\end{enumerate}
\end{enumerate}
Par définition, $a\in H^{p,q}(X)$ si et seulement s'il existe dans la classe de $a$ des formes fermées $\alpha_1$ et $\alpha_2$ dont les composantes de type $(p',q')$ sont nulles pour $p'<p$ (resp. $q'<q$). Il existe alors $\beta$ tel que $\alpha_1-\alpha_2=d\beta$. Soit $\beta_1$ (resp. $\beta_2$) la somme des composantes de $\beta$ de type $(p',q')$ pour $p'\ge p$ (resp. $q'\ge q$); on a $\beta=\beta_1+\beta_2$ et
\[
\alpha=\alpha_1-d\beta_1=\alpha_2+d\beta_2
\]
est une forme fermée de type $(p,q)$ dans la classe de $a$. Ceci prouve (i). Il suffit de prouver (ii) pour $a$ bihomogène, donc représenté par une forme bihomogène fermée $\alpha$ d'après (i). Une telle forme vérifie automatiquement $d'\alpha=d''\alpha=0$.

Si $\alpha$ vérifie $d'\alpha=d''\alpha=0$, ses composantes bihomogènes sont fermées, et chacune des conditions a), b), c) équivaut à la conjonction des conditions analogues pour les diverses composantes de $\alpha$, trivialement pour b) et c), et d'après (5.3) (ii) et (5.4) (i) pour a). Supposons donc $\alpha$ fermée de type $(p,q)$ pour prouver (iii). La classe de cohomologie $a$ définie par $\alpha$ se trouve dans $H^{p,q}(X)$, donc dans $F^p(H^{p+q}(X))$, et est nulle si et seulement si elle se trouve déjà dans $F^{p+1}(H^{p+q}(X))$, c'est-à-dire si son image dans $H^q(X,\Omega_X^p)$ est nulle. Ceci prouve a) $\Leftrightarrow$ c), et l'équivalence a) $\Leftrightarrow$ b) s'en déduit par conjugaison.

\begin{quote}\footnotesize
(Ajouté sur épreuves.) On peut montrer que ces conditions équivalent encore à l'existence d'une forme $\beta$ telle que $\alpha=d'd''\beta$.
\end{quote}

\paragraph{Théorème (5.5).} Soient $S$ un schéma de caractéristique $0$ et $f:X\to S$ un morphisme propre et lisse. Alors:
\begin{enumerate}[label=(\roman*)]
\item les faisceaux $\R^qf_*\Omega_{X/S}^p$ sont localement libres de type fini et de formation compatible à tout changement de base;
\item la suite spectrale (5.1.1)
\[
E_1^{pq}=\R^qf_*\Omega_{X/S}^p\Longrightarrow H_{\DR}^{p+q}(X/S)
\]
dégénère;
\item en chaque point de $S$, les faisceaux $\R^qf_*\Omega_{X/S}^p$ et $\R^pf_*\Omega_{X/S}^q$ ont même rang (symétrie de Hodge).
\end{enumerate}
La seconde assertion de (i) résulte de la première et de EGA III, (7.8.5) appliqué à chacun des $\Omega_{X/S}^p$; elle implique qu'il suffit de vérifier (iii) pour $S$ spectre d'un corps.

Des arguments standards permettent de se ramener successivement au cas où $S$ est affine (car la question est locale sur $S$), noethérien (par passage à la limite inductive d'anneaux), spectre d'un anneau local noethérien (car la question est locale), spectre d'un anneau local noethérien complet (par fidèle platitude de la complétion), spectre d'un anneau local artinien. Pour effectuer cette dernière réduction, on représente l'anneau local noethérien complet $R$ comme limite projective de ses quotients artiniens $R/\mathfrak m^n$. D'après le théorème de comparaison EGA III, (4.1.5), le terme $E_1$ de la suite spectrale (5.1.1) est limite projective des termes $E_1$ des suites spectrales analogues sur les $R/\mathfrak m^n$, de sorte que si celles-ci dégénèrent, (5.1.1) dégénère de même. Si l'assertion (i) est vraie sur les $R/\mathfrak m^n$, ces termes $E_1$ forment un système projectif de modules libres, se déduisant les uns des autres par réduction mod $\mathfrak m^n$, de sorte que leur limite projective est un module libre.

Supposons donc que $S=\Spec(A)$ avec $A$ anneau local artinien de corps résiduel $k$ de caractéristique $0$. Cet anneau admet un corps de représentants (EGA, $0_{IV}$, (19.8.6), (i), pour $W=k$, $u=\mathrm{identité}$) et peut donc être vu comme une $k$-algèbre locale de dimension finie. Le principe de Lefschetz permet de supposer $k=\C$, auquel cas (iii) résulte de (i) et de (5.3.3).

Pour achever la démonstration de (5.5), il reste donc à prouver la première assertion de (i), et (ii), sous l'hypothèse supplémentaire:
\[
\tag{5.5.1}
S=\Spec(A)\quad\text{pour une }\C\text{-algèbre locale de dimension finie sur }\C,
\text{ d'idéal maximal }\mathfrak m.
\]
La suite spectrale (5.1.1) se réécrit
\[
\tag{5.5.2}
E_1^{pq}=H^q(X,\Omega_{X/S}^p)\Longrightarrow \R^{p+q}\Gamma(\Omega_{X/S}^\bullet).
\]

\paragraph{Lemme (5.5.3).} Sous les hypothèses (5.5) et (5.5.1), le complexe $\Omega_{X/S}^{\an,\bullet}$ est une résolution du faisceau constant $A$ sur l'espace topologique $X^{\an}$.

Le complexe $\Omega_{X/S}^{\an,\bullet}$ étant $A$-linéaire, sa filtration $\mathfrak m$-adique est une filtration par des sous-complexes. Pour cette filtration,
\[
\Gr\Omega_{X/S}^{\an,\bullet}=\Gr A\otimes_\C \Gr^0\Omega_{X/S}^{\an,\bullet}.
\]
Le complexe $\Gr^0\Omega_{X/S}^{\an,\bullet}$ est une résolution du faisceau constant $\C$, car il n'est autre que le complexe de De Rham de $X_{\red}$. Le complexe $\Gr\Omega_{X/S}^{\an,\bullet}$ est donc une résolution de $\Gr A$, et l'assertion en résulte.

L'espace $X$ est compact, de sorte que par GAGA (cf. (5.2)) les hypercohomologies au sens algébrique et au sens transcendant de $\Omega_{X/S}^\bullet$ coïncident, et, d'après (5.5.3), coïncident encore avec la cohomologie de $X^{\an}$ à valeurs dans $A$. Dès lors,
\[
\tag{5.5.4}
\lgop_A \R^n\Gamma(\Omega_{X/S}^\bullet)=\lgop(A)\dim_\C \R^n\Gamma(\Omega_{X_{\red}}^\bullet).
\]
D'après (3.5.1) et EGA III, (6.10.5), on sait que
\[
\tag{5.5.5}
\lgop_A H^q(X,\Omega_{X/S}^p)\le
\lgop(A)\dim_\C H^q(X_{\red},\Omega_{X_{\red}}^p),
\]
l'égalité ne pouvant avoir lieu que si $H^q(X,\Omega_{X/S}^p)$ est $A$-libre.

On tire de (5.5.2) que
\[
\tag{5.5.6}
\sum_{p+q=n}\lgop_A H^q(X,\Omega_{X/S}^p)
\ge \lgop_A\R^n\Gamma(\Omega_{X/S}^\bullet),
\]
l'égalité ne pouvant avoir lieu pour tout $n$ que si la suite spectrale (5.5.2) dégénère. Mettant bout à bout ces inégalités, on trouve
\[
\tag{5.5.7}
\lgop(A)\sum_{p+q=n}\dim_\C H^q(X_{\red},\Omega_{X_{\red}}^p)
\ge
\lgop(A)\R^n\Gamma(\Omega_{X_{\red}}^\bullet),
\]
l'égalité ne pouvant avoir lieu que si les conclusions (5.4) sont vérifiées. Cette égalité résulte de (5.3) (i) appliqué à $X_{\red}$.

\section*{6. Application à la cohomologie cohérente}

\paragraph{Théorème (6.1).} Soient $S$ un schéma de caractéristique $0$ et $f:X\to S$ un morphisme projectif, lisse de dimension relative $n$. Quel que soit $p$, par exemple $p=0$, on a, dans $D^b(S,\Ocal_S)$,
\[
\R f_*\Omega_{X/S}^p\simeq \bigoplus_q \R^qf_*\Omega_{X/S}^p[-q].
\]
Désignons par $u$ la première classe de Chern d'un faisceau inversible relativement ample sur $X$, à valeurs dans $H^1(X,\Omega_{X/S}^1)$. La classe $u$ définit, par cup-produit, des homomorphismes dans $D^b(S)$,
\[
u\wedge:\R f_*\Omega_{X/S}^p\longrightarrow \R f_*\Omega_{X/S}^{p+1}[1],
\]
et un endomorphisme de degré $2$
\[
\tag{6.1.1}
u\wedge:\bigoplus_p\R f_*\Omega_{X/S}^p[-p]
\longrightarrow
\left(\bigoplus_p\R f_*\Omega_{X/S}^p[-p]\right)[2].
\]

\paragraph{Lemme (6.2).} L'endomorphisme (6.1.1) vérifie les hypothèses de (1.5).

D'après (5.4) (i) et le principe de Lefschetz, il suffit de démontrer (6.2) lorsque $S=\Spec(\C)$. L'assertion résulte alors du théorème de Lefschetz (voir (2.6.2) ou (2.6.3)) et de la décomposition de Hodge.

D'après (1.5), le complexe $\bigoplus_p\R f_*\Omega_{X/S}^p[-p]$ est isomorphe, dans la catégorie dérivée, à la somme de ses objets de cohomologie, et il en va de même pour chacun des $\R f_*\Omega_{X/S}^p$ qui, à un décalage près, en sont facteur direct (1.12).

\begin{center}\bfseries BIBLIOGRAPHIE\end{center}
\begin{enumerate}[label={\textup{[}\arabic*\textup{]}},leftmargin=2.8em]
\item Blanchard, Sur les variétés analytiques complexes, \emph{Ann. Sci. É.N.S.} \textbf{73} (1956).
\item R. Godement, \emph{Théorie des faisceaux}, Publ. Inst. Math. Univ. Strasbourg, XIII, Act. Sci. et Ind., Hermann.
\item M. Hakim, \emph{Schémas relatifs}, thèse.
\item R. Hartshorne, \emph{Residues and duality}, Lecture Notes no. 20, Springer, 1966.
\item H. Hironaka, Resolution of singularities of an algebraic variety over a field of characteristic zero, \emph{Ann. of Math.} \textbf{79} (1964).
\item S. L. Kleiman, \emph{Algebraic cycles and the Weil conjectures}, mimeographed notes, Columbia University.
\item J.-P. Serre, Géométrie algébrique et géométrie analytique, \emph{Ann. Inst. Fourier}, Grenoble, \textbf{6} (1956), cité GAGA.
\item J.-L. Verdier, \emph{Catégories dérivées}, thèse à paraître à North Holland Publ. Co.
\item J.-L. Verdier, \emph{Catégories dérivées (état 0)}, notes miméographiées, I.H.E.S.
\item A. Weil, \emph{Introduction à l'étude des variétés kählériennes}, Publ. Inst. Math. Univ. Nancago, VI, Act. Sci. et Ind., Hermann.
\item A. Weil, Sur les critères d'équivalence en géométrie algébrique, \emph{Math. Ann.} \textbf{128} (1954), p. 95--127.
\end{enumerate}
\noindent\textsc{SGA 4}. M. Artin, A. Grothendieck et J.-L. Verdier, \emph{Cohomologie étale des schémas}, Séminaire de géométrie algébrique de l'I.H.E.S. (1963--1964), à paraître à North Holland Publ. Co.

\noindent\textsc{SGA 5}. A. Grothendieck, \emph{Cohomologie $\ell$-adique et rationalité des fonctions L}, Séminaire de géométrie algébrique de l'I.H.E.S. (1964--1965).

\begin{center}\emph{Manuscrit reçu le 1\textsuperscript{er} juillet 1968.}\end{center}

\end{document}
